\noindent\textbf{Where the hypothesis fails, stated plainly.}
The design hypothesis had two parts.
First, carrying the KV cache across the switch would let the small model continue \emph{the same sentence} better than discarding or re-reading the context.
Second, molt's additions would make the simple method of~\citep{heo2026crossmodel} sufficient where it alone is not.
For both pairs we measured, it does not.

\paragraph{The transferred cache does not carry the content.}
With a mapped cache (C or D), the receiver keeps the topic and the voice of the text but loses its specifics.

\begin{center}\small
\begin{tabular}{lccc}
\toprule
 & C & D (molt) & A$'$ (re-read the text) \\
\midrule
needle, 1.5B$\to$0.5B & \res{q.15b-05b.C.needle} & \res{q.15b-05b.D.needle} & \res{q.15b-05b.A2.needle} \\
needle, 3B$\to$1.5B & \res{q.3b-15b.C.needle} & \res{q.3b-15b.D.needle} & \res{q.3b-15b.A2.needle} \\
needle at $\sim$2\,100 tokens, 3B$\to$1.5B & \res{q.3b-15b.C.longneedle} & \res{q.3b-15b.D.longneedle} & \res{q.3b-15b.A2.longneedle} \\
\bottomrule
\end{tabular}
\end{center}

molt's native anchors make the seam smaller and judge agreement higher than C's in both pairs, for example $+$\res{qp.3b-15b.D.C.agree_post}, CI \res{qp.3b-15b.D.C.agree_post.ci}.
Its task score is not distinguishable from C's (\res{qp.15b-05b.D.C.task}, CI \res{qp.15b-05b.D.C.task.ci}; \res{qp.3b-15b.D.C.task}, CI \res{qp.3b-15b.D.C.task.ci}).
``Better than C'' and ``good enough'' are different claims, and only the first one holds here.
This is not an implementation fault.
The same code path with an identity map on one model reproduces native decoding on \res{id.same} of \res{id.items} needle items (\cref{sec:heo}).
The teacher-forced gap of about a tenth of a nat per token (\cref{sec:results-intrinsic}) badly understates what is lost once the receiver generates freely.

\paragraph{Where the transfer wins, it wins at the wrong thing.}
At the prompt lengths of our main quality runs (\res{q.15b-05b.ctx.open}--\res{q.15b-05b.ctx.needle} tokens), re-prefill also resumes faster.
For 3B$\to$1.5B the times are \res{m.3b-15b.h1112.A2.resume}\,s for A$'$ and \res{m.3b-15b.h1112.D.resume}\,s for D, because D pays a second release.
At $\sim$2\,100 tokens the transfer is faster: \res{m.long.3b-15b.h1112.D.resume}\,s against \res{m.long.3b-15b.h1112.A2.resume}\,s.
On the same long context, however, it answers the needle \res{q.3b-15b.D.longneedle} where re-prefill answers \res{q.3b-15b.A2.longneedle}.
The seconds it saves are the seconds needed to get the answer right.

\paragraph{The careful release is a fix for one allocator setting, not a general requirement.}
Under MLX's default buffer cache, drop-then-load is killed in every cell of both pairs, and a plain \texttt{clear\_cache()} survives erratically or not at all.
There, synchronise-clear-wait is what makes the ordered schedule survive.
Under the 20\,MiB cache limit that mlx-swift's examples set on iOS, both naive variants survive every cell we ran.
What holds under both settings is only the \emph{order}: do not load the receiver before the sender is released.
That is a necessary condition, not a technique.

\paragraph{Nothing of molt v0.1's mapper survives.}
Its depth-linear layer pairing scores held-out NLL \res{fit.15b-05b.depth_nll} and \res{fit.3b-15b.depth_nll}, against \res{fit.15b-05b.kstar_nll} and \res{fit.3b-15b.kstar_nll} for the prior method's $R^2$-ranked top-$k$.
It scores 0 on the needle in both pairs.

\paragraph{What does hold.}
When the app's headroom is below what a co-resident switch needs (\res{m.15b-05b.need} and \res{m.3b-15b.need}\,MiB here), the prefill-hand-off shape is killed.
A release-first switch completes, mid-decode, with no retracted token, under both allocator settings and at both context lengths.
For what to put in the receiver's cache at that moment, every measurement we have favours the transcript over the transferred KV.
